\item A body of mass \(m\) was slowly hauled up the hill (Fig. 1.29) by a force \(F\) which at each point was directed along a tangent to the trajectory. Find the work performed by this force if the height of the hill is \(h\), the length of its base \(l\), and the coefficient of friction \(k\).
    \begin{center}
        \begin{tikzpicture}
            \node at (0, 0) {{hilldiagram.png}};
        \end{tikzpicture}
    \end{center}
\begin{solution}
    \begin{center}
        \begin{tikzpicture}
            \pic at (0, 0) {frame=3cm};
        \end{tikzpicture}
    \end{center}
    
    \begin{align*}
        \intertext{Let \( \mathbf{F} \), make an angle \( \theta \) with the horizontal at an arbitrary instant of time (see figure). Newton's second law in projection form along the direction of the force gives}
        F &= kmg \cos \theta + mg \sin \theta \quad \text{(because there is no acceleration of the body)}
        \intertext{As \( \mathbf{F} \parallel d\mathbf{r} \), the differential work done by the force \( \mathbf{F} \) }
        dA &= \mathbf{F} \cdot d\mathbf{r} = Fds \quad \text{(where \( ds = |d\mathbf{r}| \)}\\
        &= kmg \, ds \, (\cos \theta) + mg \, ds \, \sin \theta \\
        &= kmg \, dx + mg \, dy
        \intertext{Hence, }
        A &= kmg \int_{0}^{l} dx + mg \int_{0}^{b} dy \\
        &= kmgl + mgb = mg(kl + h) \tag{1.3}
    \end{align*}
\end{solution}
